\documentclass{scrreprt}

\usepackage{aligned-overset}
\usepackage{amsmath}
\usepackage{amsthm}
\usepackage{amssymb}
\usepackage{bm}
\usepackage[inline,shortlabels]{enumitem}
\usepackage{hyperref}
\usepackage[utf8]{inputenc}
\usepackage{multicol}
\usepackage{mathtools}
\usepackage{pdflscape}
\usepackage{physics}
\usepackage{tabularx}
\usepackage[table]{xcolor}
\usepackage{titling}
\usepackage{fancyhdr}
\usepackage{xfrac}
\usepackage{pgfplots}

\pgfplotsset{compat = newest}
\usepgfplotslibrary{fillbetween}
\usetikzlibrary{calc}


\author{Karsten Lehmann}
\date{WiSe 2025/26}
\title{Übungsblatt 01\\INF-B-120, Mathematische Methoden für Informatiker II}

\setlength{\parindent}{0pt}

\setlength{\headheight}{26pt}
\pagestyle{fancy}
\fancyhf{}
\lhead{\thetitle}
\rhead{\theauthor}
\lfoot{\thedate}
\rfoot{Seite \thepage}

\begin{document}
\paragraph{Ü 1.1}
\begin{enumerate}[(a)]
\item Bildet die Menge der ungeraden natürlichen Zahlen mit der Addition eine
  Halbgruppe?

  \subparagraph{Lsg.} Also die Menge der ungeraden natürlichen Zahlen ist schon
  mal nicht leer (1 ist zum Beispiel eine ungerade natürliche Zahl) und die
  Addition der ungeraden natürlichen Zahlen ist zwar assoziativ, allerdings nicht
  abgeschlossen, da zum Beispiel $1 + 1 = 2$ und $2$ keine ungerade natürliche
  Zahl ist.

  $\Rightarrow$ keine Halbgruppe

\item Ist $\qty(\mathbb{Z} \times \mathbb{Z}; +)$ mit komponentenweiser Addition
  eine Gruppe?

  \subparagraph{Lsg.} Es ist $\mathbb{Z} \times \mathbb{Z}$ nicht leer und
  zusammen mit der komponentenweisen Addition eine Halbgruppe, da
  die Addition auf $\mathbb{Z}$ bekannterweise assoziativ und abgeschlossen ist.

  Weiter ist $\qty\big(0, 0)^T \in \mathbb{Z} \times \mathbb{Z}$ mit
  $a + \qty\big(0, 0)^T = a$ für alle $a \in \mathbb{Z} \times \mathbb{Z}$.
  Somit ist $\qty\big(0, 0)$ neutrales Element.

  Schließlich ist $\qty(\mathbb{Z}; +)$ bereits als Gruppe bekannt, somit
  existiert für $b \in \mathbb{Z}$ auch $-b \in \mathbb{Z}$ mit
  $b + \qty\big(-b) = 0$ und schließlich für
  $\qty(b, c)^T \in \mathbb{Z} \times \mathbb{Z}$ auch
  $\qty(-b, -c)^T \in \mathbb{Z} \times \mathbb{Z}$ mit
  \[
    \qty(b, c)^T + \qty(-b, -c)^T = \qty\big(0, 0)^T
  \]
  $\Rightarrow$ Gruppe
\end{enumerate}

\paragraph{Ü 1.2} Es sei $\qty\big(G; \circ)$ die kleinste Gruppe der Ordnung
6, die die paarweise verschiedenen Elemente $a$, $b$ und $e$ enthält und für die
gilt: $e$ ist das neutrale Element der Gruppe, $a \circ a = e, b^3 = e$ und
$a \circ b = b^{-1} \circ a$.

Finden Sie alle Elemente von $G$ und stellen Sie die Verknüpfungstafel von
$\qty\big(G; \circ)$ auf.

\subparagraph{Lsg.} $G$ hat Ordnung 6, also 6 Elemente.

\begin{enumerate}[(1)]
\item Das inverse Element ist eindeutig bestimmt und die Elemente sind paarweise
  verschieden.
  Da $a$ selbstinvers ist, folgt $b^2 \ne a$
  (ansonsten $b^3 = b^2 \circ b = a \circ b = e$).
  Außerdem $b^2 \ne e$, da sonst $b^3 = b^2 \circ b = e \circ b = e$.

  Sei somit $b^2 = c$. $\Rightarrow c \circ b = b \circ c = e$.

\item Es ist $a \circ b \ne e$, da $a$ sonst zwei inverse hätte.
  Außerdem ist $a \circ b \ne c$, da sonst
  \[
    a \circ b = c \circ a \overset{b\text{ in der Mitte}}\iff
    a \circ b \circ b = c \circ b \circ a \iff
    a \circ c = e \circ a
  \]
  womit $c$ zwei inverse Elemente hätte.

  Sei daher $a \circ b = d$

\item Jetzt haben $a$, $b$, $c$ und $e$ schon bekannte Inverse.
  Somit kann $d$ entweder Selbstinvers oder Invers zu $f$ sein.

  Machen wir es invers zu $f$ um mehr Felder in der Tabelle zu fühlen.

\item Der Rest ist wie Sudoku.
\end{enumerate}

\begin{tabular}{c|cccccc}
  $\circ$ & $a$ & $b$ & $c$ & $d$ & $e$ & $f$ \\
  \hline
  $a$ & $e$ & $d$ & $f$ & $c$ & $a$ & $b$ \\
  $b$ & $f$ & $c$ & $e$ & $a$ & $b$ & $d$ \\
  $c$ & $d$ & $e$ & $b$ & $f$ & $c$ & $a$ \\
  $d$ & $c$ & $f$ & $a$ & $b$ & $d$ & $e$ \\
  $e$ & $a$ & $b$ & $c$ & $d$ & $e$ & $f$ \\
  $f$ & $b$ & $a$ & $d$ & $e$ & $f$ & $c$ \\
\end{tabular}

\paragraph{Ü 1.3} Es sind die dreielementige Menge $A \coloneqq \qty\big{a, b, c}$
und $T_A \coloneqq \qty{f {\big |} f \colon A \to A}$ gegeben.

\begin{enumerate}[(a)]
\item Zeigen Sie, dass $\qty(T_A; \circ)$ mit der Hintereinanderausführung
  $\circ$ von Abbildungen eine Halbgruppe ist.

  \subparagraph{Lsg.} $T_A$ ist nicht leer, da zum Beispiel
  $\text{id}_A \in T_A$.

  Seien nun $d, e, f \in T_A$ beliebige Abbildungen.
  Dann ist $\qty\big(d \circ e) \circ f = d \circ \qty\big(e \circ f)$, da für
  jedes $a \in A$ gilt
  \[
    \qty(\qty\big(d \circ e) \circ f)\qty\big(a)
    = \qty\big(d \circ e)\qty(f\qty\big(a))
    = d\qty(e\qty(f\qty\big(a)))
    = d\qty(\qty\big(e \circ f)\qty\big(a))
    = d \circ \qty\big(e \circ f)
  \]
  Schließlich handelt es sich ausschließlich um Abbildungen von $A$ nach $A$,
  daher bildet auch jegliche Verknüpfung der Abbildungen wieder eine Abbildung
  von $A$ nach $A$
  $\Rightarrow$ die Verknüpfung ist über $A$ abgeschlossen.

  Somit ist $\qty(T_A; \circ)$ eine Halbgruppe.

\item Zeigen Sie, dass $\qty(T_A; \circ)$ ein Monoid aber keine Gruppe ist.

  \subparagraph{Lsg.} Es ist $\text{id}_A$ das neutrale Element in
  $\qty(T_A; \circ)$, da für jedes $a \in A, f \in T_A$ gilt
  \[
    \qty\big(f \circ \text{id}_A)\qty\big(a) = f\qty(\text{id}_A\qty\big(a))
    = f\qty\big(a) =
    \text{id}_A\qty(f\qty\big(a)) = \qty(\text{id}_A \circ f)\qty\big(a)
  \]
  Somit ist $\qty\big(T_A; \circ)$ ein Monoid.

  Sei $f \in T_A$ nun die Abbildung definiert durch
  $f\qty\big(a) = f\qty\big(b) = f\qty\big(c) = a$.
  Angenommen $\qty(T_A; \circ)$ wäre eine Gruppe, dann müsste eine
  Abbildung $g \in T_A$ mit $g \circ f = \text{id}_A$ existieren.

  Somit müssten gleichzeitig
  \begin{itemize}
  \item $\qty\big(g \circ f)\qty\big(a) = g\qty(f\qty\big(a)) = g\qty\big(a) = a$ und
  \item $\qty\big(g \circ f)\qty\big(b) = g\qty(f\qty\big(b)) = g\qty\big(a) = b$
  \end{itemize}
  gelten - ein Widerspruch und somit ist $\qty(T_A; \circ)$ keine Gruppe.
\item Es werden Zwei Abbildungen $f$ und $g$, definiert durch
  \begin{flalign*}
    f&\colon A \to A \text{ mit } f\qty\big(a) = a, f\qty\big(b) = f\qty\big(c) = b \\
    g&\colon A \to A \text{ mit } g\qty\big(a) = b, g\qty\big(b) = g\qty\big(c) = a
  \end{flalign*}
  betrachtet.
  Zeigen Sie, dass $U \coloneqq \qty\big{f, g}$ eine Unterhalbgruppe von
  $\qty(T_A; \circ)$ bildet.

  \subparagraph{Lsg.} $U$ ist eine Unterhalbgruppe wenn jegliche Verknüpfung von
  Elementen aus $U$ auch wieder ein Element in $U$ bildet.
  Mit den zwei Elementen lässt sich dies schnell durchspielen:
  \begin{itemize}
  \item[$f \circ f = f$], da
    \begin{itemize}
    \item $f\qty(f\qty\big(a)) = f\qty\big(a) = a = f\qty\big(a)$
    \item $f\qty(f\qty\big(b)) = f\qty\big(b) = b = f\qty\big(b)$
    \item $f\qty(f\qty\big(c)) = f\qty\big(b) = b = f\qty\big(c)$
    \end{itemize}
  \item[$g \circ g = f$], da
    \begin{itemize}
    \item $g\qty(g\qty\big(a)) = g\qty\big(b) = a = f\qty\big(a)$
    \item $g\qty(g\qty\big(b)) = g\qty\big(a) = b = f\qty\big(b)$
    \item $g\qty(g\qty\big(c)) = g\qty\big(a) = b = f\qty\big(c)$
    \end{itemize}
  \item[$f \circ g = g$], da
    \begin{itemize}
    \item $f\qty(g\qty\big(a)) = f\qty\big(b) = b = g\qty\big(a)$
    \item $f\qty(g\qty\big(b)) = f\qty\big(a) = a = g\qty\big(b)$
    \item $f\qty(g\qty\big(c)) = f\qty\big(a) = a = g\qty\big(c)$
    \end{itemize}
  \item[$g \circ f = g$], da
    \begin{itemize}
    \item $g\qty(f\qty\big(a)) = g\qty\big(a) = b = g\qty\big(a)$
    \item $g\qty(f\qty\big(b)) = g\qty\big(b) = a = g\qty\big(b)$
    \item $g\qty(f\qty\big(c)) = g\qty\big(b) = a = g\qty\big(c)$
    \end{itemize}
  \end{itemize}
\end{enumerate}

\newpage
\paragraph{Ü 1.4} Aus der symmetrischen Gruppe $\qty(S_5; \circ)$ über der
Menge $A = \qty\big{1, 2, 3, 4, 5}$ werden die Permutationen
\[
  \alpha = \begin{pmatrix}
    1 & 2 & 3 & 4 & 5 \\
    3 & 4 & 5 & 2 & 1 \\
  \end{pmatrix}, \qquad
  \beta = \begin{pmatrix}
    1 & 2 & 3 & 4 & 5 \\
    4 & 2 & 3 & 1 & 5 \\
  \end{pmatrix}
\]
betrachtet.
\begin{enumerate}[(a)]
\item Geben Sie $\alpha, \beta, \alpha^{-1}, \alpha \circ \beta$ und
  $\beta \circ \alpha$ in Zyklenschreibweise an.

  \subparagraph{Lsg.} Es sind
  \begin{itemize}
  \item $\alpha = \qty\big(3\;5\;1)\qty\big(4\;2)$
  \item $\beta = \qty\big(4\;1)$
  \item $\alpha^{-1} = \qty\big(5\;3\;1)\qty\big(4\;2)$
  \item $\alpha \circ \beta = \qty\big(2\;4\;3\;5\;1)$
  \item $\beta \circ \alpha = \qty\big(3\;5\;4\;2\;1)$
  \end{itemize}

\item Bestimmen Sie alle Elemente der Untergruppe $\langle \alpha \rangle$
  von $\qty(S_5; \circ)$ in Zyklenschreibweise.

  Wie viele Linksnebenklassen hat $\langle \alpha \rangle$?
  Gilt
  $\beta \circ \langle \alpha \rangle = \beta^2 \circ \langle \alpha \rangle$?

  \subparagraph{Lsg.} Es sind
  \begin{itemize}
  \item $\alpha^0 = \text{id}_{S_5}$
  \item $\alpha = \qty\big(3\;5\;1)\qty\big(4\;2)$
  \item $\alpha^2 = \qty\big(3\;1\;5)$
  \item $\alpha^3 = \qty\big(4\;2)$
  \item $\alpha^4 = \qty\big(3\;5\;1)$
  \item $\alpha^5 = \qty\big(3\;1\;5)\qty\big(4\;2)$
  \end{itemize}
  und somit
  \[
    \langle \alpha \rangle = \qty{\alpha^i \:{\big |}\: i \in \mathbb{N}}
    = \qty\Big{
      \text{id}_{S_5},
      \qty\big(3\;5\;1)\qty\big(4\;2),
      \qty\big(3\;1\;5),
      \qty\big(4\;2),
      \qty\big(3\;5\;1),
      \qty\big(3\;1\;5)\qty\big(4\;2)
    }
  \]

  \textbf{Alternativ nach der Übung:} von $\alpha^1$ bis $\alpha^6$ \\

  Nun ist bekannt, dass zwei Linksnebenklassen entweder identisch oder disjunkt
  sind.
  Weiter ist $\abs{S_5} = 5! = 120$ und die Zahl der Linksnebenklassen
  $\abs{S_5} / \abs{\langle \alpha \rangle} = 120 / 6 = 20$.

\newpage
\item Bekanntlich definiert die Untergruppe $U \coloneqq \langle \alpha \rangle$
  eine Äquivalenzrelation $R_U$ in $\qty(S_5; \circ)$ durch
  \[
    \qty(\sigma, \tau) \in R_U \colon\iff \sigma^{-1} \circ \tau \in U
  \]
  \begin{enumerate}[(1)]
  \item Wie viele Elemente besitzt die Faktormenge $S_5/R_U$?

    \subparagraph{Lsg.} Die Faktormenge ist die Menge der Äquivalenzklassen
    einer Äquivalenzrelation.
    Hier entspricht eine Äquivalenzklasse jeweils einer Linksnebenklasse von
    $U$ und somit besitzt die Faktormenge 20 Elemente.

  \item Bestimmen Sie die Äquivalenzklasse
    $\qty[\qty\big(1\;5\;3)\qty\big(2\;4)]_{R_U}$.

    \subparagraph{Lsg.} Das entspricht $\alpha^5$, dementsprechend ist
    $\qty\big(1\;5\;3)\qty\big(2\;4)$ und die Äquivalenzklasse
    ist gleich $\langle \alpha \rangle$.
  \end{enumerate}
\end{enumerate}
\end{document}
